\documentclass{article}
\usepackage{amsmath, amssymb, amsthm}
\usepackage{hyperref}
\usepackage{geometry}
\geometry{margin=1in}

\title{Erd\H{o}s Problem \#1062}
\date{Last updated: 2025-09-28}
\author{Source: erdosproblems.com}

\begin{document}
\maketitle

\noindent\textbf{Status:} OPEN \quad \textbf{No prize}

\noindent\textbf{Tags:} number theory

\medskip

\noindent\textbf{Problem Statement:}

Let $f(n)$ be the size of the largest subset $A\subseteq \{1,\ldots,n\}$ such that there are no three distinct elements $a,b,c\in A$ such that $a\mid b$ and $a\mid c$. How large can $f(n)$ be? Is $\lim f(n)/n$ irrational?




The example $[m+1,3m+2]$ shows that $f(n)\geq\lceil \frac{2}{3}n\rceil$. Lebensold \cite{Le76} has shown that, for large $n$,\[0.6725 n \leq f(n) \leq 0.6736 n.\]This is problem B24 in Guy's collection \cite{Gu04}.




References

[Gu04] Guy, Richard K., Unsolved problems in number theory . (2004), xviii+437.

[Le76] Lebensold, Kenneth, A divisibility problem . Studies in Appl. Math. (1976/77), 291--294.

\noindent\textbf{Related OEIS sequences:} A038372


\bigskip
\noindent\small{Source: \url{https://www.erdosproblems.com/1062}}
\end{document}
